\section{The energy bound and the scaling contradiction}
\label{sec:conclusion}

We now close the argument. Throughout this section the reductions of
Section~\ref{sec:preliminaries} have been made, so
\(n\ge3\), \(A,B>-2\), \(pq>1\), and the gap
\(\gamma\) in Equation~\eqref{eq:gap} is positive.

\begin{proposition}[Uniform local energy bound]\label{prop:energy}
For the fixed equation parameters there is a constant \(C\), independent
of the solution pair, such that
\begin{equation}\label{eq:energy-bound}
 \int_{B_1}\bigl(|x|^B u^{q+1}+|x|^A v^{p+1}\bigr)\dd x
 \le E\le C.
\end{equation}
\end{proposition}

\begin{proof}
Write
\[
 T=\int w(x)D(x,y)\dd\pi(x,y),\qquad
 Q=\int\bigl(\Phi(x,y)-\eta(x)+w(x)D(x,y)\bigr)\dd\pi(x,y).
\]
Lemma~\ref{lem:localization} gives, for every \(\eps>0\),
\begin{equation}\label{eq:closing-errors}
 |Q|+|E_f-E_g|\le\eps E+C_\eps.
\end{equation}
The virial identity, Equation~\eqref{eq:virial}, is therefore exactly
\begin{equation}\label{eq:closing-identity}
\begin{gathered}
 c_1E_f+c_2E_g
   =aW_f+bW_g-(n-2)T+(n-2)Q,\\
 c_1=a(n+B)-(n-2),\quad c_2=b(n+A).
\end{gathered}
\end{equation}
Proposition~\ref{prop:pressure} bounds the first three terms on the
right from above by \(\eps E+C_\eps\). Consequently
\[
 c_1E_f+c_2E_g\le(n-1)\eps E+C_\eps.
\]
Here and below the additive constant may change, but remains independent
of the solution. Since \(c_1+c_2=\gamma\),
\[
 c_1E_f+c_2E_g
   =\frac{\gamma}{2}E+\frac{c_1-c_2}{2}(E_f-E_g).
\]
Using Equation~\eqref{eq:closing-errors}, we obtain
\[
 \frac{\gamma}{2}E
 \le\left(n-1+\frac{|c_1-c_2|}{2}\right)\eps E+C_\eps.
\]
Choose \(\eps\) so that the coefficient of \(E\) on the right is at
most \(\gamma/4\). This proves the universal bound on \(E\).
The first inequality in Equation~\eqref{eq:energy-bound} follows from
\(\eta=1\) on \(B_1\) and the nonnegativity of the energy densities.
\end{proof}

\begin{proof}[Proof of Theorem~\ref{thm:main}]
The elementary reductions in Section~\ref{sec:preliminaries} already
exclude \(n=2\), weights at most \(-2\), and \(pq\le1\).
Suppose a solution exists in the remaining range. For every \(R>0\),
the pair
\[
 u_R(x)=R^\alpha u(Rx),\qquad v_R(x)=R^\beta v(Rx)
\]
is an admissible solution with the same equation parameters.
Proposition~\ref{prop:energy} applies to every such pair with the same
constant. Both energy terms have the same dilation exponent:
\begin{align*}
 &\int_{B_1}
   \bigl(|x|^B u_R(x)^{q+1}+|x|^A v_R(x)^{p+1}\bigr)\dd x\\
 &\hspace{2em}
 =R^{\alpha+\beta+2-n}
   \int_{B_R}\bigl(|x|^B u(x)^{q+1}+|x|^A v(x)^{p+1}\bigr)\dd x.
\end{align*}
For \(R\ge1\), the last integral is at least its fixed, strictly
positive value on \(B_1\). It is finite on every bounded ball.
Equation~\eqref{eq:gap} and \(1-a-b>0\) give
\(\alpha+\beta+2-n>0\), so the displayed expression tends to
infinity as \(R\to\infty\). This contradicts the uniform bound.
\end{proof}

\begin{corollary}[H\'enon--Lane--Emden classification; Phan's Conjecture C]
\label{cor:phan-classification}
Let \(n\ge3\) be an integer, let \(p,q>0\), and let \(A,B>-2\).
There is no pair
\[
 u,v\in C^2(\R^n\setminus\{0\})\cap C(\R^n),
 \qquad u(x)>0,\quad v(x)>0\quad(x\in\R^n),
\]
satisfying Equation~\eqref{eq:system} at every \(x\ne0\) if and only if
\[
 \frac{n+A}{p+1}+\frac{n+B}{q+1}>n-2.
\]
When the reverse weak inequality holds, such a pair exists and can be
chosen radial.
\end{corollary}

\begin{proof}
The strict inequality gives nonexistence by Theorem~\ref{thm:main}.
Suppose instead that
\[
 \frac{n+A}{p+1}+\frac{n+B}{q+1}\le n-2.
\]
Since \(n+A,n+B>n-2>0\),
\[
 \frac1{p+1}+\frac1{q+1}
 <\frac1{n-2}\left(\frac{n+A}{p+1}+\frac{n+B}{q+1}\right)
 \le1.
\]
Thus \(pq>1\), because
\(1-1/(p+1)-1/(q+1)=(pq-1)/((p+1)(q+1))\).

We now apply Proposition~B of Phan~\cite{Phan}, which states the radial
ground-state theorem of Bidaut-V\'eron and
Giacomini~\cite[Theorem~1.4(i)]{BidautVeronGiacomini2010} for the same
equation and solution class, with \(n\ge3\), \(A,B>-2\), and \(p,q>0\).
Its radial existence region is exactly the displayed reverse weak
inequality, including equality on the critical hyperbola. In the original
Hamiltonian notation, the substitution is \(N=n\), \(a=A\), \(b=B\),
\(\delta=p\), and \(\mu=q\), with both differential-operator exponents
equal to \(2\); the preceding calculation verifies the standing
superlinearity condition \(D=\delta\mu-1>0\).

The original ground states have positive radial profiles \(U,V\) on
\((0,\infty)\), classical there and continuous at zero. This origin
regularity gives the strict positivity at zero required here: the radial
equation makes
\(r^{n-1}U'(r)\) strictly decreasing. Its limit as \(r\downarrow0\)
must be zero: a nonzero or infinite limit would contradict continuity
of \(U\) at zero after integration against \(r^{1-n}\). Hence
\(U'(r)<0\), and \(U(0)\ge U(r)>0\) for every \(r>0\). The same
argument applies to \(V\). Thus \(u(x)=U(|x|)\) and
\(v(x)=V(|x|)\) belong to the stated solution class and are strictly
positive on all of \(\R^n\). Radial existence in this class rules out
unrestricted nonexistence under the reverse weak inequality, completing
the equivalence.
\end{proof}

This classification is a consequence of Theorem~\ref{thm:main} joined
to the cited radial existence theorem, rather than a new construction
of ground states. The complementary existence assertion is restricted
to \(n\ge3\) and \(A,B>-2\); the broader nonexistence statement of
Theorem~\ref{thm:main} is unchanged.
